A different approach: Try to learn $\P^*$ directly.\\
\subtext{$\P^*$ is the data-generating distribution}

Some terminology:

\begin{tabular}{ll}
    $\mathcal{D} = \bigl\{ (x_i,y_i) \bigr\}_{i=1}^n$   & Dataset, sampled i.i.d. from $\P^*$      \\
    $\P^*$                                              & Data-generating distribution      \\
    $\mathcal{P}$                                       & Family of potential distributions \\
    $\hat{\P} \in \mathcal{P}$                          & Optimal model of $\P^*$
\end{tabular}

Some advantages \& applications:
\begin{enumerate}
    \item Allows assumptions about data-generating process\\
    \subtext{e.g. what is the likelihood of sampling $\mathcal{D}$}?
    \item Understand why some methods work\\
    \subtext{e.g. on which distributions does the square loss work?}
    \item Encode prior knowledge into the model
    \item Quantify uncertainty of predictions
    \item Develop new decision rules
    \item Generate entirely new samples
\end{enumerate}

\subsection{Assumptions}

\textbf{Assumption 1}: We assume $\mathcal{D}$ is i.i.d. sampled from $\P^*_{X,Y}$. Thus:
$$
    \P_\mathcal{D} = \prod_{i=1}^n \P_{X_i,Y_i} \qquad \text{(Independence)}
$$

\remark i.i.d. is a strong assumption: often false in practice.\\
\subtext{e.g. sampling with temporal/spatial dependencies, bias, etc.}

\method \textbf{General-purpose Estimators}\\
Methods that make no further assumptions, e.g. Histograms or Kernel Density Estimation (KDE).\\
\subtext{Generally require large $\mathcal{D}$ to be accurate, thus discouraged}

\newpage

\textbf{Assumption 2}: $\P^* \in \mathcal{P}$ (some family of param. models $\mathcal{P}$)

\definition \textbf{Parametric family of distributions}\\
\smalltext{$\theta \in \Theta \subset \R^p$ fully describes the distribution $\P^\theta$}
$$
    \mathcal{P} = \Bigl\{ \P^\theta \ \Big|\ \theta \in \Theta \Bigr\}
$$

The art here, is to choose $\mathcal{P}$ s.t. $\P^* \in \mathcal{P}$ is likely. Then:
$$
    \exists \theta^* \in \Theta:\quad \P^* = \P^{\theta^*} \in \mathcal{P}
$$
\subtext{$\theta \mapsto \P^\theta$ is assumed to be continuous. The advantage of this is that, if $\theta$ is close to $\theta^*$, then $\P^\theta$ is close to $\P^{\theta^*}=\P^*$.}

\subsection{Statistical Inference}

\textbf{Problem}: How to choose $\hat{\P}$ from $\mathcal{P}$, s.t. $\hat{\P}$ is close to $\P^*$?\\
\subtext{If $\mathcal{P}$ is parametric, this is the same as looking for $\hat{\theta} \in \Theta$ close to $\theta^*$}

{\footnotesize
    \notation if $Z$ has $\P_Z \in \mathcal{P} = \{ \P^\theta_Z \sep \theta \in \Theta \}$
    \begin{align*}
        p(z;\theta)         &= p_Z^\theta(z)
    \end{align*}
    \notation In the Bayesian context, where $\theta^*$ is sampled from $\P_\theta$:
    \begin{align*}
        p(\theta)           &= p_{\theta^*}(\theta)     \\
        p(z \sep \theta)    &= p_{Z|\theta^*=\theta}(z) \\
        p(\theta \sep z)    &= p_{\theta^*|Z=z} 
    \end{align*}
    Where $p$ is either a density or mass function.
}

There are 2 paradigms:
\begin{enumerate}
    \item \textbf{Frequentist}: Model only using observed data
    \item \textbf{Bayesian}: Model also using prior beliefs
\end{enumerate}

\subsubsection{Bayesian Paradigm}

\textbf{Further Assumption}: $\theta^*$ is sampled from a distribution $\P_{\theta^*}$\\
\subtext{Note how $\P_{\theta^*} \neq \P^{\theta^*}$.}

\theorem \textbf{Bayes' Theorem} (Applied to Inference)
$$
    \underbrace{p(\theta \sep \mathcal{D})}_\text{Posterior Belief} = \underbrace{\frac{p(\mathcal{D}\sep\theta)}{p(\mathcal{D})}}_\text{Update} \cdot \underbrace{p(\theta)}_\text{Prior Belief}
$$
$$
    p(\mathcal{D}) = \int p(\mathcal{D}\sep\theta) \cdot p(\theta)\ \text{d}\theta
$$

\subsubsection{Maximum Likelihood Estimator (MLE)}

Frequentist Approach: $\theta^*$ is considered fixed a priori.

\method \textbf{Maximum Likelihood Estimator}\\
Finds $\hat{\theta}_\text{MLE}$, which maximizes chance of observing $\mathcal{D}$ over the possible distributions $\mathcal{P} = \bigl\{ \P^\theta_{X,Y} \sep \theta \in \Theta \bigr\}$.

\definition \textbf{Maximum Likelihood Estimator}\\
\smalltext{Corresponding to $\hat{\P}_{X,Y}=\P^{\hat{\theta}_\text{MLE}}_{X,Y}$}
$$
    \hat{\theta}_\text{MLE} = \underset{\theta\in\Theta}{\text{arg max}}\ p(\mathcal{D};\theta) \overset{\text{i.i.d.}}{=} \underset{\theta\in\Theta}{\text{arg max}}\prod_{i=1}^n p\bigl( x_i,y_i;\theta \bigr) 
$$ 
{\footnotesize
    \remark Since $\log$ is strictly mon. increasing, the maximizer of the log-likelihood also maximizes the likelihood.
}

Applying several transformations:
{\footnotesize
\begin{align*}
        \hat{\theta}_\text{MLE} &= \underset{\theta\in\Theta}{\text{arg max}}\prod_{i=1}^n p\bigl( x_i,y_i;\theta \bigr)                                                                        \\
                                &= \underset{\theta\in\Theta}{\text{arg max}}\ \log \Biggl( \prod_{i=1}^n p\bigl( x_i,y_i;\theta \bigr) \Biggr)                                                 \\
                                &= \underset{\theta\in\Theta}{\text{arg max}}\ \sum_{i=1}^n \log \Bigl( p\bigl( x_i,y_i;\theta \bigr) \Bigr)                                                    \\
                                &= \underbrace{\underset{\theta\in\Theta}{\text{arg min}}\ \sum_{i=1}^n -\log \Bigl( p\bigl( x_i,y_i;\theta \bigr) \Bigr)}_\text{Generative Model}              \\
                                &= \underset{\theta\in\Theta}{\text{arg min}}\ \sum_{i=1}^n -\log \Bigl( p\bigl( y_i \sep x_i ; \theta \bigr)\cdot p(x_i) \Bigr)                                \\
                                &= \underset{\theta\in\Theta}{\text{arg min}}\ \sum_{i=1}^n -\log \Bigl( p\bigl( y_i \sep x_i ; \theta \bigr) \Bigr) + \underbrace{\sum_{i=1}^n - \log\bigl(p(x_i)\bigr)}_\text{Indep. from $\theta$}    \\
                                &= \underbrace{\underset{\theta\in\Theta}{\text{arg min}}\ \sum_{i=1}^n -\log \Bigl( p\bigl( y_i \sep x_i ; \theta \bigr) \Bigr)}_\text{Discriminative Model}
\end{align*}        
}

This has turned into an optimization problem. 2 approaches:
\begin{enumerate}
    \item Analytically: insert $p\bigl( x_i,y_i ; \theta \bigr)$ or $p\bigl( y_i \sep x_i ; \theta \bigr)$.\\
    \subtext{There are closed-form expressions in this statistical model.}
    \item Numerically: Gradient Descent
\end{enumerate}

\remark MLE is useful: It can be shown to converge to $\theta^*$.

\newpage

\subsubsection{Maximum A Poseriori Estimator (MAP)}

Bayesian Approach: $\theta^*$ is considered a random variable.

\method \textbf{Maximum A Posteriori Estimator}\\
Finds $\hat{\theta}_\text{MAP}$, which maximizes post. belief $p(\theta\sep\mathcal{D})$, i.e. it finds the $\theta \in \Theta$ with the highest density \textit{after} obataining $\mathcal{D}$.

\definition \textbf{Maximum A Posteriori Estimator}\\
\smalltext{Corresponding to $\hat{\P}_{X,Y} = \P^{\hat{\theta}_\text{MAP}}_{X,Y}$}
$$
\hat{\theta}_\text{MAP} = \underset{\theta \in \Theta}{\text{arg max}}\ p\bigl( \theta\sep\mathcal{D} \bigr)
$$
Applying several transformations:
{\footnotesize
\begin{align*}
    \hat{\theta}_\text{MAP}     &= \underset{\theta \in \Theta}{\text{arg max}}\ p\bigl( \theta\sep\mathcal{D} \bigr)                                                                                                               \\
                                &= \underset{\theta \in \Theta}{\text{arg max}}\ p\bigl( \mathcal{D}\sep\theta \bigr)\cdot p(\theta)                                                                                                \\
                                &\overset{\text{i.i.d.}}{=} \underbrace{\underset{\theta \in \Theta}{\text{arg max}}\Biggl( \prod_{i=1}^n p\bigl( x_i,y_i \sep \theta \bigr) \Biggr)\cdot p(\theta)}_\text{Generative Model}        \\ 
                                &= \underset{\theta\in\Theta}{\text{arg min}} \sum_{i=1}^n -\log\Bigl( p\bigl( x_i,y_i \sep \theta \bigr) \Bigr) - \log\bigl( p(\theta) \bigr)                                                      \\
                                &= \underset{\theta\in\Theta}{\text{arg min}} \sum_{i=1}^n -\log\Bigl( p\bigl( y_i \sep x_i , \theta \bigr) \Bigr)\cdot p\bigl( x_i \sep \theta \bigr) - \log\bigl( p(\theta) \bigr)                \\
                                &= \underset{\theta\in\Theta}{\text{arg min}} \sum_{i=1}^n -\log\Bigl( p\bigl( y_i \sep x_i , \theta \bigr) \Bigr) + \underbrace{\sum_{i=1}^n -\log\bigl( p(x_i) \bigr)}_\text{Indep. from $\theta$} - \log\bigl( p(\theta) \bigr)           \\
                                &= \underbrace{\underset{\theta\in\Theta}{\text{arg min}} \sum_{i=1}^n -\log\Bigl( p\bigl( y_i \sep x_i , \theta \bigr) \Bigr) - \log\bigl( p(\theta) \bigr)}_\text{Discriminative Model}
\end{align*}
}
{\footnotesize
    \remark Intuitively, we can use $p(\theta)$ as a weight for $\theta$, which can be used to introduce prior assumptions.
}

\lemma \textbf{MAP without prior knowledge is MLE}\\
\smalltext{Assume $\P_{\theta^*}=\mathcal{U}(\Theta)$}
{\footnotesize
$$
    \hat{\theta}_\text{MAP} = \underset{\theta\in\Theta}{\text{arg max}} \prod_{i=1}^n p\bigl( x_i,y_i \sep \theta \bigr) = \underset{\theta\in\Theta}{\text{arg max}} \prod_{i=1}^n p\bigl( x_i,y_i; \theta \bigr) = \hat{\theta}_\text{MLE}
$$
}
\subtext{Since $p(\theta)$ can thus be eliminated}

\newpage
\subsection{Bayes Optimal Predictor}
What can we do once we have $\hat{\P}$? We can estimate $\P_{\mathcal{Y}\sep X=x}$ and thus derive a decision rule $f^*(x)$.

\definition \textbf{Bayes' Optimal Predictor}\\
\smalltext{Best possible predictor when knowing $\P_{\mathcal{Y}|X}$}
$$
    f^*(x) = \underset{a \in \mathcal{Y}}{\text{arg min}}\ \E\Bigl[ l(a,\mathcal{Y}) \sep X=x \Bigr] = \underset{a\in\mathcal{Y}}{\text{arg min}}\int p\bigl( y\sep x \bigr)\cdot l(a,y)\ \text{d}y 
$$
{\footnotesize
    \remark In practice, $\P_{\mathcal{Y}|X}$ is unknown, so $\hat{\P}_{Y|X}$ is used.
}

This is the theoretically best possible predictor over all function classes $F$, an optimal solution to supervised learning:
$$
    \hat{f} = \underset{f\in F}{\text{arg min}}\sum_{i=1}^n l\Bigl( f(x_i),y \Bigr)
$$
\subtext{The proof for this is surprisingly straightforward}

\subsection{Probabilistic Perspective: Regression}